\documentclass[11pt]{article}
\usepackage[utf8]{inputenc}
\usepackage[T1]{fontenc}
\usepackage[spanish]{babel}
\usepackage[margin=2.5cm]{geometry}
\usepackage{amsmath}
\usepackage{tikz}
\usetikzlibrary{babel}

\title{Programación III \\ Tarea 6: Árboles y Grafos}
\author{Sebastián Lescano \\ 202510447}
\date{2026-2}

\begin{document}
\maketitle

Cada ejercicio está en su propio archivo (\texttt{ejercicio1.cpp} a \texttt{ejercicio4.cpp}) y se compila por separado. En todos la altura se cuenta en niveles: el árbol vacío tiene altura 0 y un solo nodo tiene altura 1.

\section{Ejercicio 1: por qué el BST se degrada a $O(n)$}

Insertar, buscar o eliminar una clave recorre un solo camino desde la raíz, así que el costo es proporcional a la altura $h$ del árbol.

Un árbol binario de altura $h$ tiene como máximo $2^h - 1$ nodos y como mínimo $h$ nodos (uno por nivel). Entonces, para $n$ nodos:
\[
\log_2(n + 1) \le h \le n
\]

El BST solo exige el orden izquierda $<$ nodo $<$ derecha, y eso no impide llegar al extremo $h = n$. Se llega insertando las claves en orden $a_1 < a_2 < \dots < a_n$. Por inducción sobre $k$:

\begin{itemize}
\item Con $k = 1$ el árbol es un solo nodo, una cadena de altura 1.
\item Si después de $k$ inserciones el árbol es una cadena donde cada nodo solo tiene hijo derecho, la clave $a_{k+1}$ es mayor que todas las anteriores, así que en cada comparación baja por la derecha y queda como hijo derecho del último nodo. El árbol sigue siendo una cadena y ahora tiene altura $k + 1$.
\end{itemize}

Por lo tanto $h = n$. La inserción número $i$ hace $i - 1$ comparaciones, y el total para construir el árbol es:
\[
\sum_{i=1}^{n} (i - 1) = \frac{n(n - 1)}{2}
\]

Eso es $O(n^2)$ en total y $O(n)$ por operación, igual que una lista enlazada. En el ejercicio 2 se ve con datos: al insertar $1, 2, \dots, 20$ en orden el BST queda con altura 20.

\section{Ejercicio 2: por qué la altura del Red-Black Tree es $O(\log n)$}

Sea $bh(x)$ la cantidad de nodos negros en cualquier camino desde $x$ hasta una hoja nula, sin contar a $x$ y contando la hoja nula. Por la propiedad 4 este número está bien definido.

\textbf{Lema.} El subárbol con raíz $x$ tiene al menos $2^{bh(x)} - 1$ nodos.

Por inducción sobre la altura de $x$:

\begin{itemize}
\item Si $x$ es una hoja nula, $bh(x) = 0$ y el subárbol tiene $2^0 - 1 = 0$ nodos.
\item Si $x$ es un nodo, cada hijo tiene altura negra $bh(x)$ si es rojo o $bh(x) - 1$ si es negro, o sea al menos $bh(x) - 1$. Por hipótesis cada subárbol hijo tiene al menos $2^{bh(x) - 1} - 1$ nodos, y sumando a $x$:
\[
2\left(2^{bh(x) - 1} - 1\right) + 1 = 2^{bh(x)} - 1
\]
\end{itemize}

Ahora se toma el camino más largo desde la raíz $r$, que tiene $h$ nodos. Por la propiedad 3 no puede haber dos rojos seguidos, y por la propiedad 2 empieza en negro, así que al menos la mitad de esos nodos son negros. Entonces $bh(r) \ge h/2$.

Juntando las dos cosas:
\begin{align*}
n &\ge 2^{bh(r)} - 1 \ge 2^{h/2} - 1 \\
n + 1 &\ge 2^{h/2} \\
h &\le 2\log_2(n + 1)
\end{align*}

Por eso $h$ es $O(\log n)$. La inserción mantiene las propiedades porque el nodo nuevo entra rojo (no cambia la altura negra) y la corrección solo recolorea o rota, y cada rotación es $O(1)$. La recoloración puede subir hasta la raíz, así que insertar cuesta $O(\log n)$.

Resultado al insertar $n = 20$ elementos en orden ascendente:

\begin{center}
\begin{tabular}{lr}
\hline
Árbol & Altura \\
\hline
BST & 20 \\
Red-Black Tree & 6 \\
Cota $2\log_2(21)$ & 8.78 \\
\hline
\end{tabular}
\end{center}

Con el recorrido pre-orden que imprime el programa se pueden dibujar los dos árboles, porque en un BST el primer valor es la raíz, los menores van a su izquierda y los mayores a su derecha.

Pre-orden del Red-Black Tree (N negro, R rojo):

\begin{center}
\small 8(N) 4(R) 2(N) 1(N) 3(N) 6(N) 5(N) 7(N) 12(R) 10(N) 9(N) 11(N) 16(N) 14(R) 13(N) 15(N) 18(R) 17(N) 19(N) 20(R)
\end{center}

\begin{figure}[ht]
\centering
\begin{tikzpicture}[
  level distance=1.1cm,
  level 1/.style={sibling distance=7cm},
  level 2/.style={sibling distance=3.5cm},
  level 3/.style={sibling distance=1.8cm},
  level 4/.style={sibling distance=1cm},
  level 5/.style={sibling distance=0.8cm},
  negro/.style={circle, draw, fill=black, text=white, minimum size=6.5mm, inner sep=0pt, font=\small},
  rojo/.style={negro, fill=red!75!black, draw=red!75!black}]
\node[negro] {8}
  child {node[rojo] {4}
    child {node[negro] {2}
      child {node[negro] {1}}
      child {node[negro] {3}}}
    child {node[negro] {6}
      child {node[negro] {5}}
      child {node[negro] {7}}}}
  child {node[rojo] {12}
    child {node[negro] {10}
      child {node[negro] {9}}
      child {node[negro] {11}}}
    child {node[negro] {16}
      child {node[rojo] {14}
        child {node[negro] {13}}
        child {node[negro] {15}}}
      child {node[rojo] {18}
        child {node[negro] {17}}
        child {node[negro] {19}
          child[missing]
          child {node[rojo] {20}}}}}};
\end{tikzpicture}
\caption{Red-Black Tree después de insertar del 1 al 20 en orden (altura 6).}
\end{figure}

Pre-orden del BST:

\begin{center}
\small 1 2 3 4 5 6 7 8 9 10 11 12 13 14 15 16 17 18 19 20
\end{center}

\begin{figure}[ht]
\centering
\begin{tikzpicture}[every node/.style={circle, draw, minimum size=6mm, inner sep=0pt, font=\scriptsize}]
\foreach \i in {1,...,20}
  \node (n\i) at (0.55*\i, -0.42*\i) {\i};
\foreach \i [evaluate=\i as \j using int(\i+1)] in {1,...,19}
  \draw (n\i) -- (n\j);
\end{tikzpicture}
\caption{BST después de insertar los mismos datos (altura 20). Cada nodo solo tiene hijo derecho.}
\end{figure}

\section{Ejercicio 3: por qué DFS es adecuado}

Un camino raíz-hoja es una rama completa del árbol, y DFS recorre el árbol justamente rama por rama: baja hasta una hoja antes de probar otra.

En la implementación la pila cumple un invariante: siempre contiene exactamente los nodos del camino desde la raíz hasta el nodo actual. Por eso:

\begin{itemize}
\item La suma se actualiza en $O(1)$: se suma al hacer push y se resta al hacer pop.
\item Cuando se llega a una hoja con suma $K$, el camino ya está en la pila y solo hay que copiarlo.
\item La memoria es $O(h)$, porque la pila nunca tiene más nodos que la altura.
\end{itemize}

Con BFS se avanzaría por niveles mezclando nodos de ramas distintas, y habría que guardar la suma y el camino de cada nodo que está en la cola, lo que ocupa $O(n)$ o más. Cada nodo entra y sale de la pila una sola vez, así que el tiempo es $O(n)$.

\section{Ejercicio 4: por qué BFS es el enfoque natural}

La cola es FIFO, así que los nodos salen en orden de profundidad: si la cola tiene todos los nodos del nivel $d$, al sacarlos se encolan todos sus hijos, que son exactamente los del nivel $d + 1$. Por inducción, antes de procesar cada nivel la cola contiene ese nivel completo y de izquierda a derecha. Las tres métricas se definen por niveles:

\begin{itemize}
\item \textbf{Árbol completo:} todos los niveles llenos menos el último, que se llena de izquierda a derecha. Eso equivale a que en el recorrido por niveles no aparezca ningún nodo después de un espacio vacío, y BFS recorre en ese mismo orden.
\item \textbf{Ancho máximo:} es el mayor número de nodos en un nivel. Por el invariante, el tamaño de la cola al empezar cada nivel es el ancho de ese nivel.
\item \textbf{Profundidad mínima:} como los nodos salen en orden no decreciente de profundidad, la primera hoja que sale de la cola es la menos profunda y se puede terminar ahí sin recorrer el resto.
\end{itemize}

Con DFS habría que recorrer todo el árbol y guardar aparte el nivel de cada nodo. Las tres funciones son $O(n)$ en tiempo y $O(w)$ en memoria, donde $w$ es el ancho máximo.

\end{document}
